\documentclass[conference]{IEEEtran}

\usepackage{iftex}
\ifPDFTeX
\usepackage[utf8]{inputenc}
\usepackage[T1]{fontenc}
\fi
\usepackage{amsmath,amssymb,bm,mathtools}
\usepackage{newtxtext,newtxmath}
\usepackage{textcomp}
\usepackage{array,booktabs,tabularx}
\usepackage{graphicx}
\usepackage{siunitx}
\usepackage{cite}
\usepackage[breaklinks=true,hidelinks]{hyperref}
\usepackage{bookmark}
\usepackage[nameinlink,capitalise,noabbrev]{cleveref}
\usepackage[final]{microtype}

\interdisplaylinepenalty=2500
\allowdisplaybreaks
\emergencystretch=2em
\sisetup{detect-weight=true,detect-family=true,per-mode=symbol}

\DeclareRobustCommand{\vct}[1]{\bm{#1}}
\DeclareRobustCommand{\mat}[1]{\bm{#1}}
\newcommand{\R}{\mathbb{R}}
\newcommand{\T}{^{\mathsf T}}
\newcommand{\Skew}[1]{\left[#1\right]_{\times}}
\newcommand{\norm}[1]{\left\lVert#1\right\rVert}
\newcommand{\diag}{\operatorname{diag}}

\title{Marine Attitude Initialization and Continuous Magnetic Self-Calibration\\
\large Mahony-Proxy Tilt Handoff, Two-Stage Magnetic Referencing, and Hard-Iron Gauge Correction (draft)}

\author{%
\IEEEauthorblockN{Mikhail S. Grushinskiy}
\IEEEauthorblockA{Independent Researcher, USA, Fair Lawn\\
Bareboat Necessities GitHub Project\\
Email: mgrouch@users.github.com}%
}

\begin{document}
\maketitle

\begin{abstract}
Marine inertial estimators are often judged by their steady-state filter equations,
yet a large fraction of persistent attitude error can be created before those
equations reach their nominal operating regime.  In waves, instantaneous specific
force is not a gravity vector, a warming attitude filter can therefore acquire a
tilt error, and a magnetic world reference learned in that tilted frame freezes the
error into the heading gauge.  A second and distinct error is produced by body-fixed
magnetometer hard iron: if startup declares the average measured horizontal field to
be north, the offset becomes part of north itself and no internally consistent yaw
filter can later recover true magnetic heading.  This paper documents the startup
and magnetic self-calibration architecture now shared by the OU--II, OU--III, and
two-frame Lie-group (TFG) marine inertial estimators in the \texttt{ocean-imu}
repository.  The applied path uses an exogenous, low-bandwidth Mahony complementary
observer to solve tilt and frame magnetic acquisition while the navigation filter is
held out of the loop; a legacy staged quaternion MEKF path remains available as a
matched ablation.  Magnetic acquisition is performed provisionally for availability
and repeated after convergence for accuracy.  A continuously running hard-iron
estimator then uses raw magnetometer data and yaw-stripped proxy tilt to solve a
regularized body-offset identification problem without adding a magnetic-bias state
to any navigation filter.  Current paired ablations use the v1.2.1 28 ft sailboat RAO records.
They separate hard-iron tracking in all three families from startup-policy
and refinement changes in TFG. Calibration-draw sweeps also record cases
where an additive correction worsens yaw. The results distinguish finite
simulation performance from the conditional observability arguments.
\end{abstract}

\begin{IEEEkeywords}
marine inertial navigation, attitude initialization, magnetometer calibration,
hard iron, Mahony filter, multiplicative extended Kalman filter, Lie-group EKF,
magnetic heading, sensor fusion, self-calibration
\end{IEEEkeywords}

\section{Introduction}
\label{sec:introduction}

An inertial attitude filter can converge after startup and still remain wrong by a
nearly constant amount.  The reason is that some startup decisions are not transient
states.  They define a reference frame.  If magnetic north is learned while the
estimated hull tilt is wrong, or if a body-fixed magnetometer offset is absorbed into
the learned world magnetic vector, the navigation filter can subsequently track that
incorrect frame with excellent innovation statistics.  Better steady-state tuning
cannot repair a gauge that the estimator itself defines as truth.

The problem is unusually acute on small marine platforms.  The accelerometer measures
specific force,
\begin{equation}
\vct f_m^{\mathcal B}
 = \mat R_{WB}\left(\vct a^{\mathcal W}-\vct g^{\mathcal W}\right)
 + \vct b_a + \vct n_a,
\label{eq:acc-model}
\end{equation}
not gravity alone.  Wave-induced translational acceleration can be comparable with the
gravity-projection error produced by several degrees of tilt.  Thus an
accelerometer-only leveling step made at one wave phase is not a reliable marine
attitude initializer.  The magnetometer has a complementary failure mode,
\begin{equation}
\vct m_i^{\mathcal B}
 = \mat R_i^{\mathsf T}\vct B^{\mathcal L}
 + \vct b_m + \vct d_i + \vct n_m,
\label{eq:mag-model}
\end{equation}
where $\vct b_m$ is body-fixed hard iron and $\vct d_i$ denotes unmodeled soft-iron
and sensor-axis distortion.  At one attitude, a change in $\vct b_m$ is
indistinguishable from a change in the unknown world field.  Motion is needed to
separate them.

The architecture studied here follows three principles.  First, solve the startup
tilt with a measurement-only observer whose correction bandwidth lies below the wave
band, rather than forcing the navigation MEKF to level itself while its translational
states are deliberately disabled.  Second, learn the magnetic world reference in a
frame built from that exogenous tilt, and relearn it after the slow observer has
converged.  Third, do not close the hard-iron identification window at startup:
accumulate the weak tilt excitation for the life of the filter and apply only the
regularized, physically plausible part of the solution.

The current implementation applies those ideas to three estimator families.  OU--II
and OU--III are quaternion multiplicative EKFs with different translational drift
regularizers.  TFG is a right-invariant two-frame Lie-group error-state EKF.  Their
attitude coordinates differ, but the startup evidence is largely outside those
coordinates because the proxy observer, magnetic acquisition, and hard-iron
identification are exogenous front-end processes.

This paper makes four contributions.  It formalizes startup as a persistent-frame
identification problem; documents the Mahony-proxy;
derives the continuous hard-iron solve and its regularization; and collects the
matched ablations that identify which errors each mechanism actually removes.

\section{Startup Is a Gauge-Establishment Problem}
\label{sec:gauge-problem}

A conventional description says that startup finds an initial quaternion.  That is
incomplete.  In the present marine estimators startup establishes at least three
quantities with different persistence:
\begin{enumerate}
  \item roll and pitch relative to gravity;
  \item a magnetic yaw gauge, i.e. which horizontal direction the filter calls north;
  \item long-lived nuisance estimates, especially accelerometer bias, that can absorb
        startup attitude error if allowed to learn too early.
\end{enumerate}
The first can be corrected continuously.  The second is dangerous because a world
magnetic reference is then used as if it were external truth.  The third is dangerous
because the residual accelerometer-bias model has a very long correlation time; a
startup mistake can therefore remain visible throughout a 20-minute record.

Let $\hat{\mat R}_{0}$ be the tilt frame used while learning a magnetic reference.
A simple startup average is
\begin{equation}
\hat{\vct B}_{0}
 = \frac{1}{N}\sum_{i=1}^{N}\hat{\mat R}_{0,i}\vct m_i.
\label{eq:reference-average}
\end{equation}
If $\hat{\mat R}_{0,i}$ has a persistent tilt error, the average rotates the measured
field into the wrong level frame.  If $\vct m_i$ also contains hard iron,
\begin{equation}
\hat{\vct B}_{0}
 = \vct B_{\rm hull}
 + \overline{\mat R}\,\vct b_m
 + \text{distortion/noise},
\label{eq:hardiron-in-reference}
\end{equation}
so the body offset is embedded directly in the vector later declared to be north.
The resulting yaw error is a gauge error, not a tracking error.

A persistent offset in the learned reference can therefore produce stable
tracking of an incorrect north direction. RMS alone cannot separate that
standing error from fluctuating tracking error.

\section{Marine Tilt Initialization}
\label{sec:tilt}

\subsection{Why instantaneous accelerometer leveling is insufficient}

For a stationary sensor, $-\vct f_m$ is an estimate of down.  In waves,
\cref{eq:acc-model} contains real orbital acceleration and the same construction
chases wave motion.  The relevant separation is therefore temporal: gravity is the
persistent direction, whereas the wave acceleration is oscillatory over the sea-state
band.  A useful startup attitude observer should correct slowly enough that it does
not rotate the hull frame to explain each wave-cycle acceleration.

\subsection{Private Mahony proxy}

The OU--II, OU--III, and TFG wrappers use a complementary attitude
observer based on the nonlinear complementary-filter structure of
Mahony--Hamel--Pflimlin \cite{Mahony2008NonlinearComplementary}.  In generic notation,
its gravity correction has the form
\begin{align}
\vct e_g &= \hat{\vct g}_{\mathcal B}\times
            \vct g_{\mathcal B,\rm meas},\\
\vct\omega_c &= \vct\omega_m-\hat{\vct b}_g+k_P\vct e_g,\\
\dot{\hat{\vct b}}_g &= -k_I\vct e_g,
\label{eq:mahony-generic}
\end{align}
with quaternion/SO(3) integration of $\vct\omega_c$.  The repository implementation
uses the conventional \texttt{two\_kp} and \texttt{two\_ki} parameterization.  The
proxy values in all three wrappers are
\begin{equation}
\texttt{two\_kp}=0.2,
\qquad
\texttt{two\_ki}=0.02.
\label{eq:proxy-gains}
\end{equation}
The proportional corner is intentionally below the wave band.  This proxy is not a
second navigation solution: it has no displacement states and is used as an
exogenous gravity-referenced frame for startup, wave-period preprocessing, and
magnetic learning.

The integral term is essential once the observer becomes an attitude initializer.
With no integral term, the implementation has the approximate static relation
\begin{equation}
\theta_{\rm bias}\simeq \frac{2b_g}{\texttt{two\_kp}}.
\label{eq:static-tilt}
\end{equation}
A dedicated test injects a constant 0.05 deg/s gyro bias: the observer settles at
0.711 deg of tilt when \texttt{two\_ki}=0 and at approximately zero with the deployed
0.02.  A vertical acceleration channel could previously tolerate that error because
its downstream processing was high-pass; an attitude seed cannot, because the tilt
is retained.

\subsection{Gravity agreement and handoff}

The proxy is not handed to the navigation filter immediately.  A gravity-agreement
gate and hold logic require a consistent tilt estimate, while the independent
wave-period/tuning front end is allowed to reach a usable operating point.  The
normal handoff is quality-driven, but a timeout is also required: a filter that never
starts in a calm or pathological record is a more severe failure than one that starts
from a conservative prior.  The OU wrappers and TFG use an earliest handoff near
8 s and a 150 s timeout, with the timeout internally prevented from truncating an
active magnetic acquisition.

At handoff, the navigation attitude covariance is deliberately anisotropic.  Proxy
tilt has finite uncertainty; yaw is either magnetically gauged or almost free.  The
OU wrappers use representative standard deviations of 0.035 rad for tilt and
0.087 rad for gauged yaw, with roughly $\pi/2$ uncertainty when no yaw reference
exists.  This prevents the first live accelerometer or magnetometer update from
fighting an unrealistically overconfident seed.

\section{Two Startup Policies}
\label{sec:policies}

\subsection{Deployed MahonyProxy path}

The wrapper-level default is \texttt{StartupInitPolicy::MahonyProxy}.  From the first
sample, the measurement-only front end runs while the navigation MEKF is untouched.
The proxy supplies tilt to the gravity gate and to magnetic accumulation.  When tilt,
magnetic north, and the tuning operating point are ready, the navigation estimator is
seeded once and enters its full live configuration in one step.  It does not spend
startup trying to estimate tilt with its translational block disabled.

This is also an exogeneity property.  The signal used to decide the magnetic reference
is not a function of the filter that will later consume that reference.  Therefore a
provisional MEKF error cannot be read out, averaged into the reference, and fed back as
an apparently self-consistent world vector.

\subsection{Optional staged quaternion MEKF path}

The original path remains implemented as
\texttt{StartupInitPolicy::StagedMekf}; it is useful as a backward-compatible mode and,
more importantly, as a matched scientific ablation.  In this mode the quaternion
MEKF (QMEKF-style multiplicative attitude filter) runs from the first IMU sample in a
deliberately degraded state: the linear navigation block is off, accelerometer-bias
learning is frozen, and accelerometer measurement noise is inflated.  The wrapper
reads that warming MEKF attitude both to decide when magnetic acquisition may start
and to define the frame in which the magnetic reference is averaged.

The failure mechanism is structural.  During the period when the linear block is off,
the attitude update has nowhere to assign real orbital acceleration except attitude
error.  A transient tilt error would be harmless if all startup products were later
re-estimated, but the magnetic gauge is persistent.  This is why the staged path can
look converged after several minutes while retaining a larger scored roll or bias
error.

The inner OU filter classes still retain staged behavior as their standalone default:
only the higher-level wrapper has enough information to perform the proxy handoff.
The deployed sea-state wrappers select \texttt{MahonyProxy}.

\section{Magnetic Reference Acquisition}
\label{sec:mag-acq}

\subsection{Yaw-stripped tilt frame}

The magnetic reference is learned in a level frame constructed from proxy roll and
pitch with yaw removed.  This has two purposes.  It prevents the proxy's unobservable
heading from contaminating the reference and makes the accumulation frame a function
of gravity alone.  For the continuous hard-iron estimator it also gives a useful
failure signal: if the vessel changes heading, the real world field moves inside the
hull-tied level frame and the model residual rises.

\subsection{Why acquisition is two-stage}

The low observer bandwidth that rejects wave acceleration also makes convergence slow.
Magnetic acquisition before tilt has settled can absorb that error into
the reference. The deployed provisional and refinement windows are fixed
implementation choices; the current ablation measures their downstream
consequences at the selected settings.

The deployed compromise is therefore:
\begin{enumerate}
  \item a \emph{provisional} magnetic reference as soon as the gravity gate permits,
        so heading becomes available before the longer refinement; and
  \item a \emph{refinement} starting at 90 s in the OU families and 30 s in TFG,
        each with a 30 s averaging window.
\end{enumerate}
Both passes are framed on the proxy, never on the navigation MEKF.  A refinement
framed on the MEKF would be self-confirming because the MEKF has already been steering
to the provisional reference that the second pass exists to replace.

\subsection{Bias-learning hold}

Accelerometer bias and tilt are weakly separable in waves.  If the bias estimator is
allowed to learn while the provisional magnetic reference is still in force, some of
the provisional tilt error is stored as accelerometer bias and persists because the
bias process is slow.  The OU wrappers therefore hold accelerometer-bias learning
until the refined reference has landed.
TFG uses the same sequencing principle, with estimator-specific timing/gating.

\section{Continuous Hard-Iron Identification}
\label{sec:hardiron}

\subsection{Identifiability from tilt excitation}

The continuous estimator uses the model implemented in
\texttt{ContinuousMagHardIronEstimator}:
\begin{equation}
\vct m_i = \mat R_i\T\vct B + \vct b + \vct n_i,
\label{eq:hi-model}
\end{equation}
where $\mat R_i$ is the yaw-stripped body-to-level tilt rotation, $\vct B$ is an
unknown field vector in that level frame, and $\vct b$ is a body-fixed hard-iron
offset.  Define exponentially weighted means
\begin{equation}
\bar{\mat A}=\overline{\mat R_i},\qquad
\bar{\vct w}=\overline{\mat R_i\vct m_i},\qquad
\bar{\vct m}=\overline{\vct m_i}.
\end{equation}
Eliminating $\vct B$ from the normal equations gives
\begin{equation}
\left(\mat I-\bar{\mat A}\T\bar{\mat A}\right)\vct b
 = \bar{\vct m}-\bar{\mat A}\T\bar{\vct w}.
\label{eq:hi-normal}
\end{equation}
Thus the magnetic field never has to become a navigation-filter state.

Let
\begin{equation}
\mat M=\mat I-\bar{\mat A}\T\bar{\mat A}.
\label{eq:M}
\end{equation}
For small roll and pitch with standard deviations $\sigma_r$ and $\sigma_p$,
\begin{equation}
\mat M\approx
\diag\!\left(\sigma_p^2,\sigma_r^2,
\sigma_r^2+\sigma_p^2\right).
\label{eq:small-angle-information}
\end{equation}
This exposes the difficulty: a hull moving only a few degrees provides an excitation
matrix of order $10^{-3}$.  A 15--30 s startup window does not make hard iron
well-conditioned.  Continuous accumulation is not merely convenient; it is the
identifiability strategy.

\subsection{Regularized solve}

The unmodeled term $\vct d_i$ in \cref{eq:mag-model} is not white sample noise.  Soft
iron and axis misalignment are modulated by tilt and therefore alias coherently into
the hard-iron solution.  Inverting a matrix of order $10^{-3}$ can amplify that model
error by orders of magnitude.  The implementation consequently solves
\begin{equation}
\hat{\vct b}
 = \left(\mat M+\lambda\mat I\right)^{-1}
   \left(\bar{\vct m}-\bar{\mat A}\T\bar{\vct w}\right),
\label{eq:ridge-solve}
\end{equation}
with
\begin{equation}
\lambda
 = \lambda_0
 + \lambda_{\rm rel}\frac{\operatorname{tr}(\mat M)}{3}.
\label{eq:ridge}
\end{equation}
The deployed values are an absolute floor $\lambda_0=4\times10^{-3}$ and relative
coefficient $\lambda_{\rm rel}=0.25$.  Scaling part of the ridge with excitation is
important: a fixed ridge shrinks the best-excited large-sea records most strongly,
which is opposite the desired behavior.

The estimator reports information
\begin{equation}
\mathcal I = N_{\rm eff}\lambda_{\min}(\mat M)
\label{eq:information}
\end{equation}
and does not offer a solution until both effective weight and information thresholds
are satisfied.  It also rejects implausibly large offsets and windows whose magnetic
model residual exceeds a configured limit.  Current defaults use a 600 s exponential
memory, minimum effective weight 500, minimum information 0.1, residual limit
3 $\mu$T, and bias magnitude below 35\% of the measured field norm.

\subsection{Continuous application}

Accumulation begins with raw magnetometer samples before startup is complete, but an
offset is not applied until the two-stage magnetic reference has finished.  This
separates identification memory from actuation.  The applied offset is slewed rather
than stepped (45 s nominal time constant in the OU implementation), reducing the
chance that a model-parameter change itself becomes an attitude transient.

Crucially, the estimator sees the \emph{raw} magnetometer, not a stream from which its
previous answer has already been subtracted.  Otherwise it would be fitting its own
correction.

\section{Correcting the Magnetic Gauge}
\label{sec:gauge-correction}

A subtle implementation choice determines whether a correct hard-iron estimate fixes
yaw at all.  The tempting update is
\begin{equation}
\vct m' = \vct m-\hat{\vct b},
\qquad
\vct B'_{\rm ref}=\vct B_{\rm ref}-\bar{\mat A}\hat{\vct b}.
\label{eq:naive-pair}
\end{equation}
Near the attitude at which the reference was learned,
$\mat R\T\bar{\mat A}\approx\mat I$, so the magnetometer innovation
\begin{equation}
\vct r_m=\vct m-\mat R\T\vct B_{\rm ref}
\end{equation}
is nearly unchanged by \cref{eq:naive-pair}.  The filter continues to agree with the
same wrong north.  A forced 10 $\mu$T offset through that path moved scored yaw by only
about 0.3 deg in the repository diagnostic: it was almost a no-op.

The deployed correction instead removes the body-fixed offset from the measurement
stream while preserving the canonical horizontal heading of the startup reference.
Only field magnitude and vertical/dip components are adjusted, and only by the delta
implied by the applied offset.  The corrected measurement then disagrees in heading
with the old gauge and the normal magnetometer update rotates yaw toward the corrected
field.  No attitude state is directly overwritten by continuous hard-iron learning.

This distinction also explains the TFG refinement operation.  Its discrete second
magnetic acquisition may apply a heading state correction through
\texttt{set\_attitude\_yaw\_absolute()}, because the refinement is explicitly
re-establishing the provisional gauge.  The continuous hard-iron path, by contrast,
changes measurement-model parameters and lets the estimator dynamics move yaw.

\section{Why the Calibration Remains Exogenous}
\label{sec:exogenous}

The hard-iron estimator reads two quantities: raw magnetometer samples and the
proxy's yaw-stripped tilt.  It reads no OU--II state, no OU--III state, and no TFG
state.  Its output changes magnetic measurement preprocessing/reference parameters;
it does not add a stochastic state, covariance block, or process model to any of the
three navigation estimators.

This matters for analysis.  The continuous calibration does not change the navigation
state dimension, $\mat F$, $\mat Q$, or the Riccati state covariance structure.  It is
an exogenous, slowly varying measurement-model parameter.  The existing local
stability arguments for the navigation dynamics therefore do not acquire an
additional estimator-state feedback loop merely because hard iron is learned.
This does not make arbitrary magnetic adaptation harmless: the residual, excitation,
plausibility, and slew gates are what bound the parameter variation that is actually
applied.

\section{OU--III Implementation}
\label{sec:ou3}

The OU--III wrapper uses \texttt{MahonyProxy}.  The proxy owns
tilt, gravity gating, and both magnetic acquisition frames.  A provisional magnetic
lock is followed by the 90 s/30 s refinement, and accelerometer-bias learning is held
until refinement completes.  The continuous hard-iron estimator then runs for the
life of the filter.

The inner OU--III MEKF does not need a magnetic-bias state.  It continues to receive a
corrected body magnetometer and a world reference.  Therefore the magnetic mechanism
is independent of OU--III's distinctive translational state
$[\vct v,\vct p,\vct S,\vct a_w]$ and of its integral-state regularizer.

\section{OU--II Implementation}
\label{sec:ou2}

OU--II now uses the same startup policy, observer gains, two-stage acquisition, and
\texttt{ContinuousMagHardIronEstimator} defaults as OU--III.  The port is deliberately
mechanical: the magnetic front end does not know whether the downstream filter
regularizes $\vct p,\vct v$ directly or regularizes an integral state $\vct S$.
The current source accumulates hard-iron statistics before startup gates, applies the
correction only after refinement, and subtracts the applied body offset before the
MEKF magnetometer measurement update.

Keeping identical magnetic defaults across the two OU families is important for fair
filter comparison.  A study that moves the hard-iron regularization or acquisition
window in one family but not the other is partly comparing magnetometer calibration,
not translational estimation.

\section{TFG Implementation}
\label{sec:tfg}

TFG uses different attitude-error coordinates but now follows the same front-end
logic.  Its previous magnetic acquisition consisted effectively of a single noisy
sample once the gravity gate passed.  The current wrapper uses the Mahony proxy,
windowed magnetic averaging, two-stage acquisition, and continuous hard-iron tracking.
The proxy remains outside the Lie-group state.

At refinement TFG must distinguish a heading state correction from a world-frame gauge
transformation.  The implementation uses
\texttt{set\_attitude\_yaw\_absolute()} to correct heading while leaving the world
motion columns $\vct v,\vct p,\vct S,\vct a_w$ untouched.  Rotating those columns as
though the entire world frame had changed would alter displacement history in response
to a correction that says only that the estimated hull heading was wrong.

TFG also exposes the same continuous hard-iron configuration and drives it from raw
magnetometer plus proxy tilt.  As in the OU filters, correction is applied only after
the refined magnetic reference is in force.

\section{Experimental Evidence and Attribution}
\label{sec:results}

% Generated by tools/startup_publication_sync.py.
All current ablations use the pinned v1.2.1 vessel-response records and the
final 900 s. Each row below reports the arithmetic mean of eight per-record
RMS scores; ``worst'' is the largest record RMS. These means differ from
root-mean-square pooling used by the engine study.

\subsection{Continuous hard-iron correction}
\begin{table}[t]\centering
\caption{Default-draw hard-iron ablation with current deployed startup.}
\label{tab:startup-hardiron}\footnotesize
\begin{tabular}{@{}lrrrr@{}}\toprule
& \multicolumn{2}{c}{Mean yaw RMS [deg]} & \multicolumn{2}{c}{Worst yaw RMS [deg]} \\
Family & off & on & off & on \\
\midrule
OU--II & 2.137 & 0.581 & 2.198 & 0.727 \\
OU--III & 1.945 & 0.478 & 2.064 & 0.815 \\
TFG & 1.733 & 0.531 & 1.846 & 0.926 \\
\bottomrule\end{tabular}\end{table}

\subsection{TFG feature attribution}
TFG retains both startup policies, allowing matched feature-off replays.
The OU wrappers retain only their deployed proxy startup, so no numerical
comparison to a removed OU startup implementation is reported.
\begin{table}[t]\centering
\caption{TFG default-draw feature ablations. Bias is mean accelerometer-bias RMS.}
\label{tab:tfg-ablation}\footnotesize
\resizebox{\columnwidth}{!}{%
\begin{tabular}{@{}lrrrrr@{}}\toprule
Arm & Yaw [deg] & Roll [deg] & Pitch [deg] & 3-D [m] & Bias [m/s$^2$] \\
\midrule
Deployed & 0.5313 & 0.3085 & 0.0683 & 0.3038 & 0.0516 \\
Hard iron off & 1.7330 & 0.2984 & 0.0666 & 0.3036 & 0.0497 \\
Mag refinement off & 0.7591 & 0.4652 & 0.2905 & 0.3041 & 0.0956 \\
Staged startup & 2.0054 & 0.2696 & 0.0966 & 0.3066 & 0.0484 \\
\bottomrule\end{tabular}}\end{table}

\subsection{Calibration-draw sensitivity}
The five draws are default, 11, 23, 202 and 3033. Each changes the initial
magnetic offset and distortion while holding the paired realization fixed.
For OU--II, mean yaw RMS changes from 3.025 to 1.769 degrees; 3 of 40 pairs worsen.
For OU--III, mean yaw RMS changes from 2.933 to 1.735 degrees; 4 of 40 pairs worsen.
These finite draws do not prove uniform calibration improvement.


\section{What Improved, and Why}
\label{sec:what-improved}

The paired feature-off arms identify the effect of each retained switch
under the declared sensor draws. The generated tables report all attitude,
bias and displacement channels; a single yaw score does not describe the
full cost of a startup choice. These observations do not identify the effect
of retired OU startup implementations.

\section{Robustness and Failure Modes}
\label{sec:limits}

\subsection{Soft iron and sensor misalignment}

The continuous estimator identifies only an additive body-fixed offset.  It cannot
separate that offset perfectly from multiplicative distortion when heading does not
change. The current five-draw measurements are reported above, including
all pairs with worse yaw. An offset-only model cannot identify
soft-iron/misalignment error without additional heading excitation.
A misalignment-dominated residual may alias into the additive solution.

The failure is not reliably detectable from the hard-iron residual because the
aliased component is precisely the part the simplified model can explain.  This is
why ridge shrinkage is part of the model-error bound rather than merely numerical
regularization.

\subsection{Heading maneuvers}

The current accumulation frame is tied to the hull's heading.  A real turn moves the
world magnetic vector inside that frame and raises the residual; the estimator
therefore stands down.  A future estimator could exploit heading excitation to
identify a full soft-iron/misalignment model, but that is a different observability
problem and is not claimed here.

\subsection{Magnetic disturbances}

Local ferromagnetic changes, current-dependent magnetic fields, and nearby equipment
can violate the assumption of a slowly varying body-fixed hard-iron vector.  Residual,
field-norm, plausibility, and slew gates are necessary but cannot make an arbitrary
magnetic environment identifiable.  Operational implementations should expose
calibration validity and avoid interpreting every yaw innovation as evidence for a
new permanent offset.

\subsection{Very weak wave excitation}

The hard-iron solve needs tilt diversity.  In a nearly motionless vessel,
$\lambda_{\min}(\mat M)$ becomes small and the information gate prevents the estimator
from claiming a precise offset.  This is preferable to inverting an unobservable
problem.  The long memory allows later motion to contribute information without
requiring the device to restart calibration.

\section{Recommended Validation Contract}
\label{sec:validation}

A publication-quality extension of this work should preserve feature attribution.
At minimum, future evaluations should include:
\begin{enumerate}
  \item paired \texttt{MahonyProxy} versus \texttt{StagedMekf} startup with all later
        magnetic options held fixed;
  \item magnetic refinement on/off with startup and hard iron fixed;
  \item continuous hard iron on/off with startup and refinement fixed;
  \item multiple IMU noise seeds and independent magnetometer calibration draws;
  \item fixed-heading wave records and deliberate heading maneuvers;
  \item additive hard iron, multiplicative soft iron, axis misalignment, and mixed
        calibration errors as separate oracle-controlled arms;
  \item startup latency, time to first valid heading, handoff attitude covariance,
        and recovery after a large tilt reset;
  \item per-axis roll/pitch/yaw RMS, accelerometer-bias error, and displacement metrics
        so gains cannot be hidden in a single aggregate attitude number.
\end{enumerate}
The TFG one-feature-off study is a useful model because it demonstrates that the
largest yaw improvement and the largest roll/bias improvement come from different
mechanisms.

\section{Discussion}
\label{sec:discussion}

The main design lesson is that initialization and calibration should be assigned to
the estimator that has the right information, not automatically to the estimator
that ultimately consumes the result.  The navigation MEKF has the richer state, but
during startup that richness is intentionally disabled or poorly observed.  A small
independent complementary observer can therefore be the better source of tilt for a
one-time magnetic decision because it is not entangled with the provisional magnetic
reference or with translational states that are not yet live.

The second lesson is that a magnetometer world reference is not just another tuning
constant.  Its horizontal direction defines a gauge.  Once a body-fixed offset has
been absorbed into that direction, the filter can be statistically self-consistent
and physically wrong.  The continuous hard-iron estimator works because it solves the
sensor-identification problem outside the navigation filter and then changes the
measurement model in a way that forces a new heading innovation instead of preserving
the old one.

The third lesson is that continuous self-calibration does not necessarily require
augmenting the navigation state.  When a nuisance parameter can be identified from
exogenous sufficient statistics, keeping it outside the navigation covariance can
preserve estimator structure, simplify cross-family reuse, and make ablation much
cleaner.  That strategy is justified here specifically for a slowly varying hard-iron
offset under the stated excitation and distortion limits; it is not a general argument
against joint-state calibration.

\section{Conclusion}
\label{sec:conclusion}

Marine attitude startup is a persistent-frame estimation problem, not merely a request
for an initial quaternion.  Wave acceleration corrupts instantaneous gravity leveling;
a magnetic reference learned in the resulting tilt frame can freeze that error into
roll, pitch, yaw, and slowly learned accelerometer bias.  The deployed OU--II,
OU--III, and TFG architecture separates the problem into an exogenous low-bandwidth
Mahony proxy, two-stage magnetic acquisition, delayed nuisance-bias learning, and
continuous regularized hard-iron identification.  The legacy staged quaternion MEKF
path remains valuable as a matched ablation and demonstrates why the separation is
needed.

The measurements show distinct benefits.  Proxy startup primarily reduces roll/tilt
and accelerometer-bias error without materially moving OU displacement.  Continuous
hard-iron correction then removes the standing magnetic gauge: mean yaw falls by
about 55--57\% in both OU families, and TFG's hard-iron-off arm more than doubles yaw
while leaving the other channels nearly unchanged.  The remaining limitation is also
clear: additive hard-iron learning cannot identify soft-iron/misalignment error without
additional heading excitation.  That boundary is preferable to hiding the problem
inside a larger filter.  The resulting architecture is therefore best understood as
staged observability management: gravity first establishes tilt, a refined magnetic
reference establishes heading, and long-horizon motion supplies the weak excitation
needed to continuously improve the sensor calibration.

\bibliographystyle{IEEEtran}
\bibliography{w3d}

\end{document}
